\documentclass[10pt,a4paper]{amsart}
\usepackage[margin=19mm]{geometry}
\usepackage{amsmath,amssymb,mathtools}
\usepackage[scr=rsfs,cal=euler]{mathalfa}
\usepackage{xfrac}
\usepackage[unicode,hidelinks,bookmarks=false]{hyperref}
\newcommand{\ie}{i.e.\ }
\newcommand{\cf}{cf.\ }
\newcommand{\resp}{respectively\ }
\newcommand{\Subsection}{\subsection}
\providecommand{\nobreakdash}{\nobreak\hbox{-}}
\setlength{\emergencystretch}{2em}
\multlinegap0pt
\usepackage{amssymb}
\newtheorem{theorem}{Theorem}
\newtheorem{lemma}{Lemma}
\newcommand{\prob}{\mathbb{P}}
\title{On the number of finite additive 2-bases}
\author{Stefan Weltge, Konrad Zyhalko}
\date{Source: arXiv:2605.19449v2 (CC BY 4.0)}
\begin{document}
\maketitle

\noindent\emph{Source and licence.} On the number of finite additive 2-bases. Version arXiv:2605.19449v2 (CC BY 4.0), distributed by arXiv under the Creative Commons Attribution 4.0 International licence, http://creativecommons.org/licenses/by/4.0/, as recorded in the arXiv licence metadata for this exact version and shown on the abstract page https://arxiv.org/abs/2605.19449v2. Full licence notice: \texttt{licenses/arxiv-2605.19449-CC-BY-4.0.txt}.\par\smallskip

\noindent\textit{Editorial scope.} Counted unit: the article's theorem thmMain (source0.tex lines 48--51). The article's complete original proof of it is delivered with it (source0.tex lines 143--191, one \texttt{proof} environment of 49 lines, which the article places away from the statement). No mathematics is written, repaired or re-derived: the statement and the proof are the article's own lines, in the article's own notation, and no retained line is substituted or re-flowed.  Retained uncounted as the source-local context the proof needs: lines 60--66; lines 97--100; lines 123--129.  Nothing was omitted from any retained span, so the package carries no omission record.  No retained line contains a citation, so the wrapper carries no bibliography and no bibliography entry.  No retained line contains a figure environment, an image-inclusion command or drawing code, so no image is delivered and the package declares no asset dependency.  Editorial formatting only, in addition: an AMS \texttt{amsart} preamble that restates the article's own declarations and macros used by the retained lines, namely the environments \texttt{theorem}, \texttt{lemma} on separate counters, together with the standard packages those lines need.  The retained environments print in this document as Theorem 1, Lemma 1 in order of appearance, whereas the article prints the same items with the numbering of its own full document.  The complete native source of the article as distributed in the arXiv e-print for this exact version is bundled unchanged as \texttt{sources/200916/source0.tex}.
\begin{lemma}
    \label{lemNotGenerated}
    If $X$ is a uniformly random subset of $[n]_0$ and $k \in [n + 1]_0$, then
    \[
        \prob[k \notin X + X] \le \left( \frac{3}{4} \right)^{\frac{k}{2} - 1}.
    \]
\end{lemma}

\begin{lemma}
    \label{lemFirstLowerBound}
    We have $|\Gamma(n)| \ge 2^n / (n+1)^{20}$.
\end{lemma}

\begin{lemma}
    \label{lemRecurrence}
    For every integer $n \ge 0$, we have
    \[
        |\Gamma(n+1)| = |\Gamma(n)| + |\{X\in \Gamma(n) : n + 1 \in X+X\}|.
    \]
\end{lemma}

\begin{theorem}
    \label{thmMain}
    There exists some $\alpha > 0$ such that $|\Gamma(n)| \ge \alpha \cdot 2^{n+1}$ holds for all integers $n \ge 0$.
\end{theorem}

\begin{proof}[Proof of Theorem~\ref{thmMain}]
    Consider the ratio
    \[
        \delta(n) := \frac{|\{X \in \Gamma(n) : n + 1 \in X + X\}|}{|\Gamma(n)|}.
    \]
    By Lemma~\ref{lemRecurrence}, we have
    \begin{equation}
        \label{eqRecurrenceDelta}
        |\Gamma(n+1)| = (1 + \delta(n)) \cdot |\Gamma(n)|.
    \end{equation}
    By Lemma~\ref{lemNotGenerated} and Lemma~\ref{lemFirstLowerBound}, we have
    \begin{align*}
        1 - \delta(n)
        = \frac{|\{X \in \Gamma(n) : n + 1 \notin X + X\}|}{|\Gamma(n)|}
        & \le \frac{|\{X \subseteq [n]_0 : n + 1 \notin X + X\}|}{|\Gamma(n)|} \\
        & \le 2 \cdot \left( \frac{3}{4} \right)^{\frac{n-1}{2}} \cdot (n+1)^{20},
    \end{align*}
    and hence
    \begin{equation}
        \label{eqDeltaBound}
        1 + \delta(n) \ge 2 - \underbrace{\tfrac{4}{\sqrt{3}} \left( \tfrac{\sqrt{3}}{2} \right)^n \cdot (n+1)^{20}}_{=: t(n)}.
    \end{equation}
    Note that
    \[
        \frac{t(n+1)}{t(n)} = \frac{\sqrt{3}}{2} \cdot \left( 1 + \frac{1}{n+1} \right)^{20},
    \]
    and hence there exists some integer $n_0$ such that
    \begin{equation}
        \label{eqLe1}
        t(n) \le \frac{1}{10}
    \end{equation}
    and
    \begin{equation}
        \label{eqGeometricDecay}
        t(n + 1) \le \frac{9}{10} t(n)
    \end{equation}
    holds for all $n \ge n_0$.
    We conclude
    \begin{align*}
        \frac{|\Gamma(n)|}{2^n}
        = \frac{|\Gamma(n_0)|}{2^{n}} \cdot \prod_{k=n_0}^{n - 1} (1 + \delta(k))
        & \ge \frac{|\Gamma(n_0)|}{2^{n_0}} \cdot \prod_{k=n_0}^{n - 1} \left( 1 - \frac{t(k)}{2} \right) \\
        & \ge \frac{|\Gamma(n_0)|}{2^{n_0}} \cdot \left( 1 - \frac{1}{2} \sum_{k=n_0}^{n - 1} t(k) \right) \\
        & \ge \frac{|\Gamma(n_0)|}{2^{n_0}} \cdot \left( 1 - \frac{t(n_0)}{2} \sum_{k=0}^{n - n_0 - 1} \left( \frac{9}{10}\right)^k \right) \\
        & \ge \frac{|\Gamma(n_0)|}{2^{n_0}} \cdot \left( 1 - 5 \cdot t(n_0) \right) \\
        & \ge \frac{|\Gamma(n_0)|}{2^{n_0 + 1}},
    \end{align*}    
    where the first equality follows from \eqref{eqRecurrenceDelta}, the first inequality from~\eqref{eqDeltaBound} and~\eqref{eqLe1}, the third inequality from~\eqref{eqGeometricDecay}, and the last inequality from~\eqref{eqLe1}.
\end{proof}

\end{document}
